\chapter{Glossário Matemático e Técnico}
\label{apendice:glossario_matematico}

Este apêndice reúne definições formais dos principais termos matemáticos,
estatísticos e de aprendizado de máquina empregados ao longo do livro.
Cada entrada inclui notação, descrição e, quando pertinente, referência
bibliográfica. Os termos aparecem em ordem alfabética.

% ---------------------------------------------------------------
\section*{B}
% ---------------------------------------------------------------

\subsection*{Byte-Pair Encoding (BPE)}
\label{sec:bpe}
\index{Byte-Pair Encoding|textbf}

O \textit{Byte-Pair Encoding} (BPE) é um algoritmo de tokenização por
sub-palavras originalmente proposto para compressão de dados e posteriormente
adaptado ao processamento de linguagem natural. O procedimento inicia com um
vocabulário de caracteres individuais e, iterativamente, funde o par de
símbolos adjacentes de maior frequência no corpus de treinamento, criando um
novo símbolo composto. Após $k$ fusões, obtém-se um vocabulário de tamanho
$|\mathcal{V}| = |\mathcal{V}_0| + k$, em que $\mathcal{V}_0$ é o vocabulário
inicial de caracteres.

\begin{enumerate}
  \item Inicializar $\mathcal{V}$ com todos os caracteres únicos do corpus.
  \item Contar a frequência de todos os pares de símbolos adjacentes.
  \item Fundir o par mais frequente $(a, b)$ em um novo símbolo $ab$.
  \item Repetir os passos 2–3 até atingir $|\mathcal{V}|$ desejado.
\end{enumerate}

O BPE oferece equilíbrio entre vocabulários de caracteres (sem palavras
desconhecidas, mas sequências longas) e vocabulários de palavras inteiras
(sequências curtas, mas alta taxa de \textit{out-of-vocabulary}).

% ---------------------------------------------------------------
\section*{E}
% ---------------------------------------------------------------

\subsection*{Embedding (vetor de representação)}
\label{sec:embedding}
\index{embedding|textbf}

Um \textit{embedding} é uma função $f: \mathcal{X} \to \mathbb{R}^d$ que
mapeia elementos de um espaço discreto (tokens, frases, documentos) para
vetores densos em $\mathbb{R}^d$. No contexto de LLMs, a camada de
\textit{embedding} de entrada $\mathbf{E} \in \mathbb{R}^{|\mathcal{V}| \times d}$
associa cada token $t \in \mathcal{V}$ a um vetor $\mathbf{e}_t \in \mathbb{R}^d$.
A dimensão $d$ varia tipicamente entre 768 (modelos base) e 4\,096 (modelos
de grande porte). Geometricamente, tokens semanticamente próximos tendem a
ocupar regiões próximas no espaço vetorial, propriedade explorada pela busca
por similaridade de cosseno:
\[
  \text{sim}(\mathbf{u}, \mathbf{v}) = \frac{\mathbf{u} \cdot \mathbf{v}}
  {\|\mathbf{u}\| \, \|\mathbf{v}\|}.
\]

% ---------------------------------------------------------------
\section*{L}
% ---------------------------------------------------------------

\subsection*{Logits}
\label{sec:logits}
\index{logits|textbf}

O termo \textit{logits} designa os valores reais não normalizados $\mathbf{z} \in \mathbb{R}^{|\mathcal{V}|}$
produzidos pela última camada linear do modelo antes da aplicação da função
de ativação. Formalmente, dado o vetor de estado oculto $\mathbf{h}$ da última
camada do \textit{Transformer}, os logits são calculados como:
\[
  \mathbf{z} = \mathbf{W}_\text{lm}\,\mathbf{h} + \mathbf{b},
\]
em que $\mathbf{W}_\text{lm} \in \mathbb{R}^{|\mathcal{V}| \times d}$ é a
matriz de projeção do modelo de linguagem e $\mathbf{b}$ é o viés. Os logits
são convertidos em probabilidades pela função \textit{softmax}
(ver Seção~\ref{sec:softmax}).

% ---------------------------------------------------------------
\section*{S}
% ---------------------------------------------------------------

\subsection*{SentencePiece}
\label{sec:sentencepiece}
\index{SentencePiece|textbf}

O \textit{SentencePiece} é uma biblioteca de tokenização por sub-palavras
independente de idioma desenvolvida pelo Google. Ao contrário do BPE clássico
(ver Seção~\ref{sec:bpe}), o \textit{SentencePiece} opera diretamente sobre
o fluxo de caracteres Unicode sem pré-segmentação por espaços, o que o torna
aplicável a idiomas sem delimitadores explícitos de palavras (japonês, chinês,
tailandês). O algoritmo suporta dois modelos: BPE e \textit{unigram language
model}, que seleciona o vocabulário minimizando a perda de verossimilhança do
corpus sob um modelo de linguagem de unigrama.

\subsection*{Softmax}
\label{sec:softmax}
\index{softmax|textbf}

A função \textit{softmax} $\sigma: \mathbb{R}^n \to \Delta^{n-1}$ converte
um vetor de logits $\mathbf{z} = (z_1, \ldots, z_n)$ em uma distribuição de
probabilidade sobre $n$ classes:
\[
  \sigma(\mathbf{z})_i = \frac{e^{z_i}}{\sum_{j=1}^{n} e^{z_j}}, \quad i = 1, \ldots, n.
\]
No contexto da geração de texto, $n = |\mathcal{V}|$ e cada $\sigma(\mathbf{z})_i$
representa a probabilidade do token $i$ ser o próximo token gerado. A
\emph{temperatura} $\tau > 0$ escala os logits antes da \textit{softmax}:
\[
  \sigma(\mathbf{z}/\tau)_i = \frac{e^{z_i/\tau}}{\sum_j e^{z_j/\tau}}.
\]
Para $\tau \to 0$, a distribuição concentra-se no token de maior logit
(\textit{greedy decoding}); para $\tau > 1$, a distribuição torna-se mais
uniforme, aumentando a diversidade das saídas.

% ---------------------------------------------------------------
\section*{W}
% ---------------------------------------------------------------

\subsection*{WordPiece}
\label{sec:wordpiece}
\index{WordPiece|textbf}

O \textit{WordPiece} é um algoritmo de tokenização por sub-palavras
desenvolvido pelo Google e utilizado em modelos como BERT. Em vez de fundir
o par de maior frequência absoluta (como no BPE), o \textit{WordPiece}
seleciona o par $(a, b)$ que maximiza o ganho de verossimilhança do modelo
de linguagem ao substituir a sequência $ab$ por um único token, definido como:
\[
  \text{score}(a, b) = \frac{\text{freq}(ab)}{\text{freq}(a) \cdot \text{freq}(b)}.
\]
Essa formulação favorece fusões de pares cujos componentes individuais são
relativamente raros, produzindo um vocabulário de sub-palavras com maior
coerência linguística.
